\chapter{Blokčejn kao infrastrukturni sloj distribuiranih sistema}

\section{Kriptografske i sistemske osnove blokčejn tehnologije}
Digitalni dokumenti u tekstualnom, audio, slikovnom i video formatu podložni su lakoj i često neprimetnoj izmeni, što nameće problem pouzdanog utvrđivanja vremena njihovog nastanka ili poslednje modifikacije. Ovaj izazov verifikabilnog vremenskog obeležavanja digitalnih podataka u nepouzdanom okruženju rešen je kroz koncepte kriptografskog vremenskog obeležavanja, pri čemu se sadržaj dokumenata najpre kriptografski hešira bez njegovog otkrivanja, a zatim povezuje sa prethodnim zapisima putem heš vrednosti, čime se formira lanac kriptografski zavisnih vremenskih oznaka \cite{haber1990time, bayer1993improving, massias2002timestamp}. Ovakva struktura omogućava detekciju svake naknadne izmene, budući da bi promena jednog zapisa narušila konzistentnost svih narednih veza u lancu. Ključni doprinos ovih radova ogleda se u demonstraciji da se integritet i hronološki redosled podataka mogu garantovati oslanjanjem na kriptografske mehanizme i javno proverljive dokaze, bez potrebe za potpunim poverenjem u centralni autoritet.

Masovno slanje elektronske pošte bez značajnog troška omogućilo je ranu zloupotrebu otvorenih računarskih sistema kroz nepoželjne poruke \eng{spam}. Kao odgovor na ovaj problem, predloženi su mehanizmi u kojima se izvršavanje operacija od značaja uslovljava rešavanjem računski zahtevnih, ali lako proverljivih zadataka. Ovaj pristup uvodi ekonomsku barijeru pristupu resursima, stvarajući asimetriju u kojoj legitimni korisnici snose zanemarljiv dodatni trošak, dok napadači koji pokušavaju da generišu veliki broj zahteva moraju da plate proporcionalno visoku cenu u vidu procesorskog vremena i potrošnje energije \cite{Dwork1992PricingVP, back2002hashcash}. Ovi radovi formalizuju koncept dokazivanja uloženog rada kao mehanizma za regulaciju ponašanja u distribuiranim sistemima i pokazuju da se bezbednost sistema može postići korišćenjem ekonomskog troška, bez oslanjanja na identitet učesnika ili centralnu kontrolu.

Pokušaji realizacije elektronskog novca sa svojstvima fizičkog gotovinskog novca, kao što su privatnost i nepratljivost transakcija, zasnivali su se na upotrebi slepih potpisa i kriptografskih protokola koji omogućavaju izvršavanje plaćanja bez otkrivanja identiteta korisnika. Međutim, ovi sistemi suočili su se sa fundamentalnim problemom dvostruke potrošnje \eng{double spending}, odnosno mogućnošću da se isti digitalni token iskoristi više puta \cite{chaum1983blind}. Predložena rešenja uglavnom su se oslanjala na centralizovane entitete koji vode evidenciju o već potrošenim jedinicama, čime se implicitno uvodi potreba za poverenjem u treću stranu \cite{chaum1985security}. Kasniji predlozi \cite{szabo2005bit, dai1998b} nastojali su da redukuju ulogu centralnog autoriteta, ali nisu u potpunosti razrešili problem koordinacije i postizanja globalne saglasnosti u otvorenom distribuiranom okruženju. Ovi radovi jasno identifikuju dvostruku potrošnju kao centralni tehnički izazov u dizajnu digitalnih monetarnih sistema.
